\documentclass[a4paper,12pt]{article}
\usepackage[utf8]{inputenc}
\usepackage{amsmath}
\usepackage{graphicx}
\usepackage{float}
\usepackage{hyperref}
\usepackage{amsfonts}
\usepackage{amssymb}
\usepackage{geometry}
\usepackage{listings}
\usepackage{xcolor}
\usepackage{cancel}
% Pacote para códigos em Python
\usepackage{listings}
\renewcommand{\lstlistingname}{Algoritmo}
% Configuração do ambiente listings para Python
\lstset{
    language=Python,
    basicstyle=\ttfamily\small,
    keywordstyle=\color{blue}\bfseries,
    stringstyle=\color{red},
    commentstyle=\color{gray}\itshape,
    numbers=left,
    numberstyle=\tiny\color{gray},
    stepnumber=1,
    numbersep=5pt,
    backgroundcolor=\color{white},
    frame=single,
    tabsize=4,
    showspaces=false,
    showstringspaces=false,
    breaklines=true,
    breakatwhitespace=true,
    captionpos=b
}
\pagestyle{empty}

\geometry{a4paper, margin=0.5in}

\title{Gabarito da Prova 2 - Álgebra Linear para Computação}
\date{}
\newcommand{\avg}{\mathbf{\textbf{avg}}}
\newcommand{\std}{\mathbf{\textbf{std}}}
\newcommand{\vet}[1]{\mathbf{#1}}
\begin{document}
\maketitle
\vspace{-4em}
\begin{enumerate}
    \item Podemos escrever a decomposição da seguinte forma:
    $$
    \begin{bmatrix}
        a_{11} & a_{12} \\
        a_{21} & a_{22} \\
    \end{bmatrix}=
    \begin{bmatrix}
        \ell_{11} & 0 \\
        \ell_{21} & \ell_{22} \\
    \end{bmatrix}
    \begin{bmatrix}
        u_{11} & u_{12} \\
        0 & u_{22} \\
    \end{bmatrix}
    $$
    
    Sabemos que a digonal de $L$ é formada de 1, portanto, temos $\ell_{11} = \ell_{22} = 1$.
    $$
    \begin{bmatrix}
        a_{11} & a_{12} \\
        a_{21} & a_{22} \\
    \end{bmatrix}=
    \begin{bmatrix}
        1 & 0 \\
        \ell_{21} & 1 \\
    \end{bmatrix}
    \begin{bmatrix}
        u_{11} & u_{12} \\
        0 & u_{22} \\
    \end{bmatrix}
    $$
    Escrevendo as multiplicações das linhas de $L$ pelas colunas de $U$, temos:
    
    \begin{align*}
        &u_{11} = a_{11}\\
        &u_{12} = a_{12}\\
        &\ell_{21} u_{11} = a_{21} \Rightarrow \ell_{21} = \frac{a_{21}}{a_{11}}\\
        &\ell_{21} u_{12} + u_{22} = a_{22} \Rightarrow u_{22} = a_{22} -  \frac{a_{21} a_{12}}{a_{11}}\\
    \end{align*}
    Aplicando essas equações na matriz dada, temos:

    $$
    \begin{bmatrix}
        5 & 2 \\
        3 & 1  \\
    \end{bmatrix} = 
        \begin{bmatrix}
        1 & 0 \\
        \frac{3}{5} & 1 \\
    \end{bmatrix}
    \begin{bmatrix}
        5 & 2 \\
        0 & -\frac{1}{5} \\
    \end{bmatrix}.
    $$
    Sabendo que $\det(A) = \det(LU) = \det(L)\det(U)$ e que a determinante de uma matriz trianguilar é o produto dos elementos da diagonal, temos $\det(A) = (1 \cdot 1) \cdot (5 \cdot \frac{-1}{5}) = -1$.


    \item 
    \begin{enumerate}
        \item Como $A$ é uma matriz simétrica, o teorema espectral nos garante a existencia das matrizes $P$ e $D$ onde $D$ é uma matriz diagonal onde seus elementos correspondem aos autovalores da matriz $A$ e as colunas de $P$ são os autovetores correspondentes. Desta forma, vamos calcular os autovalores de $A$ a partir das raízes do polinômio característico:
    \begin{align*}
        \det(A - \lambda I) &= (6-\lambda) \cdot (3-\lambda) - 4 = 0\\
        &\Rightarrow 14 - 9\lambda + \lambda^2 = 0 \\
        &\Rightarrow \lambda = \frac{9 \pm \sqrt{25}}{2} \\
        &\Rightarrow \lambda_1 = 7,\ \lambda_2 = 2. \\
    \end{align*}
    Para calcular os autovetores, temos:
    \begin{align*}
    &\begin{bmatrix}
        6 & 2\\
        2 & 3\\
    \end{bmatrix} \begin{bmatrix}
        x_1\\x_2
    \end{bmatrix} = \lambda \begin{bmatrix}
        x_1\\x_2
    \end{bmatrix}\\
    \Rightarrow & \begin{cases}
        6 x_1 + 2 x_2 = \lambda x_1 \Rightarrow x_1 = \frac{-2 }{6-\lambda}x_2\\
        2 x_1 + 3 x_2 = \lambda x_2  \Rightarrow x_2 = \frac{-2 }{3-\lambda}x_1 \\
    \end{cases}
    \end{align*}
    Para $\lambda_1$ e considerando $x_1$ como variável livre, temos: $x_2 = \frac{1}{2}x_1 \Rightarrow v_1 = \begin{bmatrix}
        1 \\ 0.5
    \end{bmatrix}$.

    Para $\lambda_2$ e considerando $x_1$ como variável livre, temos: $x_2 = -2 x_1 \Rightarrow v_2 = \begin{bmatrix}
        1 \\ -2
    \end{bmatrix}$.
    
    Para construir $P$ de forma que seja uma matriz unitária, devemos normalizar os vetores $v_1$ e $v_2$:
    \begin{align*}
        p_1 = \frac{v_1}{\|v_1\|} = \frac{v_1}{\sqrt{1 + 0.25}} = \frac{2}{\sqrt{5}} v_1 = \frac{1}{\sqrt{5}}\begin{bmatrix}
            2 \\1
        \end{bmatrix}\\
        p_2 = \frac{v_2}{\|v_2\|} = \frac{v_2}{\sqrt{1 + 4}} = \frac{1}{\sqrt{5}} v_2 = \frac{1}{\sqrt{5}}\begin{bmatrix}
            1 \\-2
        \end{bmatrix}
    \end{align*}
    Portanto, $P = \frac{1}{\sqrt{5}}\begin{bmatrix}
        2 & 1\\
        1 & -2\\
    \end{bmatrix}$
    \item $P P^T = \frac{1}{5}\begin{bmatrix}
        2\cdot2 +1 \cdot 1 & 2\cdot1 +1 \cdot (-2)\\
        1\cdot2 -2 \cdot 1 & 1\cdot1 + (-2) \cdot (-2)\\
    \end{bmatrix} = \begin{bmatrix}
        1 & 0\\
       0 & 1\\
    \end{bmatrix}$, ou seja, $P$ é uma matriz unitária.
    \item $\det(A) = \lambda_1 \cdot \lambda_2 = 7 \cdot 2 = 14$ e $\text{Tr}(A) = \lambda_1 + \lambda_2 = 7 + 2 = 9$
    \end{enumerate}
    
    \item A matriz $D$ é uma matriz diagonal onde suas entradas são os autovalores da matriz $A$. A matriz $P$ é uma matriz inversível onde suas colunas são os autovetores da matriz $A$ e a ordenação depende da matriz $D$.
    
    Vamos calcular os autovalores de $A$ a partir das raízes do polinômio característico:
    \begin{align*}
        \det(A - \lambda I) &= (1-\lambda) \cdot (3-\lambda) - 8 = 0\\
        &\Rightarrow -5 - 4\lambda + \lambda^2 = 0 \\
        &\Rightarrow \lambda = \frac{4 \pm \sqrt{36}}{2} \\
        &\Rightarrow \lambda_1 = 5,\ \lambda_2 = -1. \\
    \end{align*}

    Para calcular os autovetores, temos:
    \begin{align*}
    &\begin{bmatrix}
        1 & 4\\
        2 & 3\\
    \end{bmatrix} \begin{bmatrix}
        x_1\\x_2
    \end{bmatrix} = \lambda \begin{bmatrix}
        x_1\\x_2
    \end{bmatrix}\\
    \Rightarrow & \begin{cases}
        x_1 + 4 x_2 = \lambda x_1 \Rightarrow x_1 = \frac{-4 }{1-\lambda}x_2\\
        2 x_1 + 3 x_2 = \lambda x_2  \Rightarrow x_2 = \frac{-2 }{3-\lambda}x_1 \\
    \end{cases}
    \end{align*}
    Para $\lambda_1$ e considerando $x_1$ como variável livre, temos: $x_2 = \frac{-2}{3-5} x_1= x_1 \Rightarrow p_1 = \begin{bmatrix}
        1 \\ 1
    \end{bmatrix}$.

    Para $\lambda_2$ e considerando $x_2$ como variável livre, temos: $x_1 =\frac{-4 }{1+1} x_2= -2 x_2 \Rightarrow p_2 = \begin{bmatrix}
        -2 \\ 1
    \end{bmatrix}$.

    Portanto, $P = \begin{bmatrix}
        1 & -2 \\
        1 & 1
    \end{bmatrix}$ e, utilizando a fórmula da inversa de matrizes $2\times 2$, temos $P^{-1} = \frac{1}{\det(P)} \begin{bmatrix}
        1 & 2\\
        -1 & 1
    \end{bmatrix} = \frac{1}{3} \begin{bmatrix}
        1 & 2\\
        -1 & 1
    \end{bmatrix}$. Dessa forma, temos a autodecomposição dada por:

    $$
    \begin{bmatrix}
        1 & 4\\
        1 & 3
    \end{bmatrix} = \begin{bmatrix}
        1 & -2 \\
        1 & 1
    \end{bmatrix}\ \begin{bmatrix}
        5 &  0\\
        0 & -1
    \end{bmatrix}\ \begin{bmatrix}
         \frac{1}{3} &  \frac{2}{3}\\
         \frac{-1}{3} &  \frac{1}{3}
    \end{bmatrix}.
    $$

    Para calcular $A^4$, observe que:
    \begin{align*}
        A^4 &= (P D P^{-1})^4 \\
        &= (P D P^{-1})(P D P^{-1})(P D P^{-1})(P D P^{-1}) \\
        &= P D \underbrace{P^{-1}P}_{I} D \underbrace{P^{-1}P}_{I} D \underbrace{P^{-1}P}_{I} D P^{-1} \\
        &= P D^4 P^{-1}
    \end{align*}
    
    Portanto:
    \begin{align*}
        A^4 &= \begin{bmatrix}
        1 & -2 \\
        1 & 1
    \end{bmatrix}\ \begin{bmatrix}
        5^4 &  0\\
        0 & (-1)^4
    \end{bmatrix}\ \begin{bmatrix}
         \frac{1}{3} &  \frac{2}{3}\\
         \frac{-1}{3} &  \frac{1}{3}
    \end{bmatrix}\\
    &=\begin{bmatrix}
        1 & -2 \\
        1 & 1
    \end{bmatrix}\ \begin{bmatrix}
        625 &  0\\
        0 & 1
    \end{bmatrix}\ \begin{bmatrix}
         \frac{1}{3} &  \frac{2}{3}\\
         \frac{-1}{3} &  \frac{1}{3}
    \end{bmatrix}\\
    &=\begin{bmatrix}
        1 & -2 \\
        1 & 1
    \end{bmatrix}\ \begin{bmatrix}
         \frac{625}{3} &  \frac{1250}{3}\\
         \frac{-1}{3} &  \frac{1}{3}
    \end{bmatrix}\\
    &=\frac{1}{3} \begin{bmatrix}
         627 &  1248\\
         624 &  1251
    \end{bmatrix} = \begin{bmatrix}
         209 &  416\\
         208 &  417
    \end{bmatrix}\\
    \end{align*}

    \item Sabemos que a solução dos mínimos quadrados é dada por $x = (\vet{a}^T\vet{a})^{-1}\vet{a}^T\vet{b} = \frac{\vet{a}^T\vet{b}}{\|\vet{a}\|^2}$. Dessa forma, temos:
    \begin{align*}
        \|\vet{a}x - \vet{b}\|^2 &= \|x\vet{a}\|^2 + \|\vet{b}\|^2 -2 (x\vet{a}^T\vet{b})\\
        &=\left(\frac{\vet{a}^T\vet{b}}{\|\vet{a}\|^2} \right)^2 \|\vet{a}\|^2 + \|\vet{b}\|^2 -2 \frac{\vet{a}^T\vet{b}}{\|\vet{a}\|^2}(\vet{a}^T\vet{b})\\
        &=\frac{(\vet{a}^T\vet{b})^2}{(\|\vet{a}\|^2)^{\cancel{2}}}\  \cancel{ \|\vet{a}\|^2} + \|\vet{b}\|^2 -2 \frac{\vet{a}^T\vet{b}}{\|\vet{a}\|^2}(\vet{a}^T\vet{b})\\
        &=\cancel{\frac{(\vet{a}^T\vet{b})^2}{\|\vet{a}\|^2}} + \|\vet{b}\|^2 -\cancel{2} \frac{(\vet{a}^T\vet{b})^2}{\|\vet{a}\|^2}\\
        &= \|\vet{b}\|^2 -\frac{(\vet{a}^T\vet{b})^2}{\|\vet{a}\|^2}\\
        &= \|\vet{b}\|^2 \left(1 -\frac{(\vet{a}^T\vet{b})^2}{\|\vet{a}\|^2\|\vet{b}\|^2}\right)\\
    \end{align*}
    Repare que $\cos(\theta) = \frac{\vet{a}^T\vet{b}}{\|\vet{a}\|\|\vet{b}\|} \Rightarrow (\cos(\theta))^2 = \frac{(\vet{a}^T\vet{b})^2}{\|\vet{a}\|^2\|\vet{b}\|^2}$. Além disso, $(\cos(\theta))^2 + (\sin(\theta))^2 = 1 \Rightarrow( \sin(\theta))^2 = 1 - (\cos(\theta))^2$. Portanto:
    \begin{align*}
        \|\vet{a}x - \vet{b}\|^2 &=  \|\vet{b}\|^2 \left(1 -\frac{(\vet{a}^T\vet{b})^2}{\|\vet{a}\|^2\|\vet{b}\|^2}\right)\\
        &= \|\vet{b}\|^2 \left(1 -(\cos(\theta))^2\right)\\
        &= \|\vet{b}\|^2 (\sin(\theta))^2\\
    \end{align*}

    \item
    \begin{enumerate}
        \item A função objetivo dos mínimos quadrados ponderados pode ser escrita como:
        $$
        \sum_{i=1}^{m} w_i (\tilde{a}_i^T x - b_i)^2 = \sum_{i=1}^{m} (\sqrt{w_i}(\tilde{a}_i^T x -b_i))^2 
        $$
        Isso é equivalente a $\|D(Ax - b)\|^2$ onde

        $$
        D = \begin{bmatrix}
            \sqrt{w_1} &  0 & \cdots &0 \\
             0& \sqrt{w_2} & \cdots &0 \\
              \vdots& \vdots & \ddots & \vdots \\
            0 &0 & \cdots & \sqrt{w_m}\\
        \end{bmatrix}.
        $$
        \item Reescrevendo o problema como $\|B x - d\|^2$ onde $B = DA$ e $d = Db$, a solução do problema de mínimos quadrados é dada por $\hat{x} = (B^TB)^{-1}B^Td$. Substituindo por $A$ e $b$, temos:
        \begin{align*}
            \hat{x} &=  (B^TB)^{-1}B^Td\\
            &= ((DA)^TDA)^{-1}(DA)^TDb\\
            &= (A^TD^TDA)^{-1}A^TD^TDb\\
            &= (A^TD^2A)^{-1}A^TD^2b\\
            &= (A^TWA)^{-1}A^TWb\\
        \end{align*}

        onde $W = D^2 = \text{diag}(w) =  \begin{bmatrix}
            {w_1} &  0 & \cdots &0 \\
             0& {w_2} & \cdots &0 \\
              \vdots& \vdots & \ddots & \vdots \\
            0 &0 & \cdots & {w_m}\\
        \end{bmatrix}.$
    \end{enumerate}
\end{enumerate}

\end{document}